\documentclass[11pt,a4paper]{article}

% ── Codificación y fuentes ────────────────────────────────────────────────────
\usepackage[utf8]{inputenc}
\usepackage[T1]{fontenc}
\usepackage{lmodern}
\usepackage[spanish]{babel}

% ── Geometría y espaciado ─────────────────────────────────────────────────────
\usepackage[top=2.5cm, bottom=2.5cm, left=2.8cm, right=2.8cm]{geometry}
\usepackage{setspace}
\setstretch{1.15}
\usepackage{parskip}

% ── Tablas ────────────────────────────────────────────────────────────────────
\usepackage{booktabs}
\usepackage{tabularx}
\usepackage{array}
\usepackage{longtable}
\usepackage{enumitem}

% ── Colores y cajas ───────────────────────────────────────────────────────────
\usepackage[dvipsnames]{xcolor}
\usepackage{tcolorbox}
\tcbuselibrary{skins, breakable}

% ── Cabeceras y pies de página ────────────────────────────────────────────────
\usepackage{fancyhdr}
\usepackage{lastpage}
\pagestyle{fancy}
\fancyhf{}
\fancyhead[L]{\small\textbf{Informe de Evaluación} --- Física para Bioanalistas}
\fancyhead[R]{\small Yilselys Naranjo}
\fancyfoot[C]{\small Página \thepage\ de \pageref{LastPage}}
\renewcommand{\headrulewidth}{0.4pt}

% ── Hiperenlaces ──────────────────────────────────────────────────────────────
\usepackage[colorlinks=true, linkcolor=NavyBlue, urlcolor=NavyBlue]{hyperref}

% ── Paleta de colores personalizados ─────────────────────────────────────────
\definecolor{desempeno_excelente}{HTML}{1A6B2F}
\definecolor{desempeno_bueno}{HTML}{2E5A9C}
\definecolor{desempeno_suficiente}{HTML}{8B6914}
\definecolor{desempeno_deficiente}{HTML}{C0392B}
\definecolor{desempeno_insuficiente}{HTML}{7B241C}
\definecolor{grisFondo}{HTML}{F5F5F5}
\definecolor{azulTitulo}{HTML}{1B2A6B}
\definecolor{verdeFortaleza}{HTML}{1A6B2F}
\definecolor{naranjaMejora}{HTML}{A04000}

% ── Estilos de cajas ─────────────────────────────────────────────────────────
\tcbset{
  cajaTitulo/.style={
    enhanced, breakable,
    colback=azulTitulo!8, colframe=azulTitulo,
    fonttitle=\bfseries\large, coltitle=white,
    attach boxed title to top left={yshift=-2mm, xshift=4mm},
    boxed title style={colback=azulTitulo},
    top=6pt, bottom=6pt, left=8pt, right=8pt,
  },
  cajaFortaleza/.style={
    enhanced, breakable,
    colback=verdeFortaleza!6, colframe=verdeFortaleza!60,
    fonttitle=\bfseries, coltitle=verdeFortaleza,
    top=4pt, bottom=4pt, left=8pt, right=8pt,
  },
  cajaMejora/.style={
    enhanced, breakable,
    colback=naranjaMejora!6, colframe=naranjaMejora!60,
    fonttitle=\bfseries, coltitle=naranjaMejora,
    top=4pt, bottom=4pt, left=8pt, right=8pt,
  },
  cajaRecomendacion/.style={
    enhanced, breakable,
    colback=grisFondo, colframe=gray!50,
    fonttitle=\bfseries, coltitle=black,
    top=4pt, bottom=4pt, left=8pt, right=8pt,
  },
}

% ─────────────────────────────────────────────────────────────────────────────
\begin{document}

% ══ PORTADA ══════════════════════════════════════════════════════════════════
\begin{center}
  {\Large\bfseries\color{azulTitulo} Universidad de Oriente/Núcleo Sucre --- Física para Bioanalistas}\\[4pt]
  {\large Informe de Evaluación Preliminar}\\[12pt]
  \rule{\linewidth}{1.2pt}\\[6pt]
  {\LARGE\bfseries Yilselys Naranjo}\\[4pt]
  \rule{\linewidth}{0.6pt}
\end{center}

% ══ FICHA TÉCNICA ════════════════════════════════════════════════════════════
\vspace{4pt}
\begin{tcolorbox}[cajaTitulo, title=Datos del informe]
\begin{tabular}{@{}llll@{}}
  \textbf{Estudiante:}  & Yilselys Naranjo  &
  \textbf{Fecha:}       & 2026-08-08 \\[4pt]
  \textbf{Archivo PDF:} & \texttt{Yilselys\_Naranjo.pdf}  &
  \textbf{Páginas:}     & 6 \\
\end{tabular}
\end{tcolorbox}

% ══ RESULTADO GLOBAL ═════════════════════════════════════════════════════════
\vspace{6pt}
\begin{center}
  \begin{tcolorbox}[
    enhanced,
    colback=white, colframe=desempeno_suficiente,
    width=0.62\linewidth,
    boxrule=2pt,
    halign=center, valign=center,
    top=8pt, bottom=8pt,
  ]
    \centering
    {\large\bfseries Puntaje sugerido}\\[6pt]
    {\Huge\bfseries\color{desempeno_suficiente} 6.7/10}\\[6pt]
    {\large Nivel de desempeño: \textbf{\color{desempeno_suficiente} Suficiente}}
  \end{tcolorbox}
\end{center}

% ══ RESUMEN GENERAL ══════════════════════════════════════════════════════════
\section*{Resumen general}
La estudiante demuestra un excelente dominio procedimental, un manejo impecable de las unidades físicas y una sólida capacidad para resolver cálculos matemáticos complejos (como se evidencia en las preguntas 1, 2, 3 y 5). Sin embargo, el puntaje global se ve severamente afectado debido a que dejó una sección en blanco (2a) y, de manera crítica, respondió a una versión completamente diferente de la Pregunta 4 (desarrollando un problema sobre glucosa al 1\% en lugar de agua destilada y solución salina de NaCl), lo que sugiere una posible memorización de un banco de preguntas o confusión de examen.

% ══ FORTALEZAS Y ÁREAS DE MEJORA ════════════════════════════════════════════
\vspace{4pt}
\begin{minipage}[t]{0.48\linewidth}
  \begin{tcolorbox}[cajaFortaleza, title={Fortalezas}]
\begin{itemize}[leftmargin=1.5em, itemsep=2pt]
  \item Excelente precisión en los cálculos numéricos y uso correcto de la notación científica.
  \item Impecable manejo de unidades en el Sistema Internacional y conversiones (ej. de Joules a Kilocalorías en Q1, y metros a centímetros en Q5).
  \item Comprensión clara de la Primera Ley de Fick y su análisis de proporcionalidad (Q3a).
  \item Excelente análisis físico del fenómeno de capilaridad y la Ley de Jurin (Q5).
\end{itemize}
  \end{tcolorbox}
\end{minipage}
\hfill
\begin{minipage}[t]{0.48\linewidth}
  \begin{tcolorbox}[cajaMejora, title={Areas de mejora}]
\begin{itemize}[leftmargin=1.5em, itemsep=2pt]
  \item Atención rigurosa a las instrucciones y enunciados del examen para evitar responder preguntas de otras versiones.
  \item Evitar dejar preguntas en blanco; siempre es preferible plantear la hipótesis o el modelo físico aplicable (Q2a).
  \item Profundizar en las justificaciones físicas cualitativas, evitando explicaciones circulares (como en Q1a).
\end{itemize}
  \end{tcolorbox}
\end{minipage}

% ══ OBSERVACIONES POR PREGUNTA ═══════════════════════════════════════════════
\section*{Observaciones por pregunta}

\begin{longtable}{>{\bfseries}p{2.8cm} p{1.6cm} p{7.0cm} p{2.8cm}}
  \toprule
  \textbf{Pregunta} & \textbf{Puntaje} & \textbf{Evaluación} & \textbf{Errores detectados} \\
  \midrule
  \endhead
  \bottomrule
  \endfoot
    \textbf{Pregunta 1: Termorregulación y Eficiencia de la Sudoración} & 1.3/2.0 & En el inciso (a), la conclusión es correcta (Paciente B), pero la justificación es circular y no menciona el concepto clave de calor latente de vaporización ($L_v$) ni la diferencia entre evaporación y goteo. En el inciso (b), identifica correctamente la dependencia de la presión de vapor, pero no detalla el peligro clínico (hipertermia/golpe de calor). El cálculo del inciso (c) es excelente y preciso en ambas unidades. & \textit{Justificación física incompleta y circular en el inciso a.; Falta detallar el riesgo clínico en el inciso b.} \\
    \midrule
    \textbf{Pregunta 2: Mecánica Respiratoria y Ley de Laplace} & 1.3/2.0 & El inciso (a) fue dejado completamente en blanco, perdiendo el puntaje de la demostración de la Ley de Laplace. El inciso (b) presenta una excelente y detallada explicación del papel del surfactante. El cálculo del inciso (c) es totalmente correcto y con las unidades adecuadas (Pascales). & \textit{Inciso a no respondido (en blanco).} \\
    \midrule
    \textbf{Pregunta 3: Optimización en Hemodiálisis y Primera Ley de Fick} & 1.8/2.0 & El análisis matemático del inciso (a) es impecable, demostrando con claridad el aumento del 40\% en la tasa de transferencia. En el inciso (b), identifica correctamente la pérdida de resistencia estructural, aunque se pudo profundizar más en el riesgo clínico (ruptura y mezcla de fluidos). El cálculo del inciso (c) es correcto. & \textit{Falta de profundidad en el análisis de riesgo clínico del inciso b.} \\
    \midrule
    \textbf{Pregunta 4: Errores de Infusión Intravenosa y Ósmosis} & 0.3/2.0 & El estudiante desarrolló una pregunta completamente diferente a la planteada en el examen impreso. El examen solicita analizar agua destilada, NaCl al 5\% y calcular la presión osmótica de NaCl al 0.9\%. El estudiante respondió analizando una solución de glucosa al 1\%, calculando su molaridad y su presión osmótica. Se otorga un puntaje mínimo de 0.3 únicamente porque en el inciso (a) asocia correctamente el concepto de hipotonicidad con la entrada de agua y la lisis celular. & \textit{Resolución de un problema ajeno al examen (glucosa al 1\% en lugar de agua destilada y NaCl).; No responde a los efectos del NaCl al 5\% (crenación).; No calcula la presión osmótica del NaCl al 0.9\%.} \\
    \midrule
    \textbf{Pregunta 5: Capilaridad y Toma de Muestras Sanguíneas} & 2.0/2.0 & Respuesta perfecta. En el inciso (a) analiza correctamente la relación entre cohesión/adhesión y el ángulo de contacto para predecir la depresión capilar. En el inciso (b) explica correctamente la relación inversa entre el radio y la altura. El cálculo del inciso (c) es impecable. & \textit{---} \\
    \midrule
\end{longtable}

% ══ ERRORES DETECTADOS ═══════════════════════════════════════════════════════
\section*{Análisis de errores}

\subsection*{Errores conceptuales}
\textit{Ninguno identificado.}

\subsection*{Errores procedimentales}
\begin{itemize}[leftmargin=1.5em, itemsep=2pt]
  \item Omitir por completo el desarrollo del inciso 2a.
  \item Uso de datos y fórmulas correspondientes a otro enunciado en la Pregunta 4.
\end{itemize}

% ══ RECOMENDACIONES AL ESTUDIANTE ════════════════════════════════════════════
\section*{Retroalimentación para el estudiante}
\begin{tcolorbox}[cajaRecomendacion, title={Dirigido a: Yilselys Naranjo}]
Yilselys, tienes una excelente capacidad para el desarrollo matemático, el despeje de fórmulas y el uso de unidades físicas, lo cual es una gran fortaleza en física médica. Sin embargo, debes tener mucho cuidado al leer los enunciados del examen. En la Pregunta 4 respondiste a un problema de glucosa al 1\% que no correspondía a este examen, lo que afectó gravemente tu calificación. Asimismo, evita dejar preguntas en blanco (como la 2a); siempre intenta plantear las ecuaciones fundamentales. Sigue practicando la argumentación física cualitativa para que tus justificaciones sean tan fuertes como tus cálculos.
\end{tcolorbox}

% ══ NOTA PARA EL PROFESOR ════════════════════════════════════════════════════
\section*{Nota para el profesor}
\begin{tcolorbox}[
  enhanced, breakable,
  colback=yellow!8, colframe=yellow!60!black,
  fonttitle=\bfseries, coltitle=black,
  title={Observaciones para revisión manual},
  top=4pt, bottom=4pt, left=8pt, right=8pt,
]
Atención: En la Pregunta 4, la estudiante responde con absoluta precisión a un problema de infusión de glucosa al 1\% (calculando su presión osmótica en 1.40 atm), el cual NO es el que aparece en la hoja del examen (que pide agua destilada y NaCl 0.9\%). Esto sugiere que la estudiante memorizó la resolución de un examen de otro semestre o de otra versión y la transcribió sin leer el enunciado actual. Se recomienda una revisión manual y, de ser necesario, una breve entrevista con la estudiante para verificar la autoría de sus respuestas.
\end{tcolorbox}

% ══ PIE DE INFORME ════════════════════════════════════════════════════════════
\vspace{12pt}
\begin{center}
  \footnotesize\color{gray}
  Informe generado automáticamente por el sistema de evaluación con IA (gemini-3.5-flash) y soporte LaTex. \\
  Este documento es una evaluación \textbf{preliminar} para ser revisado por el profesor antes de ser entregado al 
  estudiante. \\
  Generado el 2026-08-08 \textbullet\ BIOANALISIS Sistema Académico de Evaluación.
  \end{center}

\end{document}
